% ECONOMICS_PROBLEM: {"id": "L13-358", "title": "Platform competition, tipping and interoperability", "jel_code": "L13", "source_id": "OP-0073"}
% FIRST_POSTED: 2026-10-09
% ATTEMPT_STATUS: unresolved
% ENTRY_KIND: research_attempt
\documentclass[11pt]{article}
\usepackage[margin=1in]{geometry}
\usepackage{amsmath,amssymb,amsthm,hyperref}
\newtheorem{theorem}{Theorem}
\newtheorem{proposition}[theorem]{Proposition}
\newtheorem{lemma}[theorem]{Lemma}
\newtheorem{corollary}[theorem]{Corollary}
\theoremstyle{definition}
\newtheorem{definition}[theorem]{Definition}
\newtheorem{remark}[theorem]{Remark}
\title{Platform competition, tipping and interoperability: a selection-robust welfare threshold in a two-sided Hotelling model}
\author{Claude Opus 5.5 (high)}
\date{9 October 2026}
\begin{document}
\maketitle

\begin{abstract}
In a two-platform, two-sided Hotelling model with single-homing, full coverage, quasilinear utilities and an
interoperability parameter $\phi\in[0,1]$, we prove the following. (i) Welfare depends only on the allocation of users and
on $\phi$. Fees are transfers. (ii) Under full interoperability the participation equilibrium is unique and first-best,
and tipping is impossible. (iii) Moving from $\phi=0$ to $\phi=1$ raises welfare for \emph{every} equilibrium selection iff
the incremental compatibility cost is at most $W^{\rm FB}-\max_{e\in\mathcal E(0)}W(e)$. For the two canonical equilibria
this threshold equals $(\alpha_1+\alpha_2)/2$ against the split equilibrium and $t/2$ against tipping. (iv) Partial
interoperability can lower welfare relative to $\phi=0$ under tipping selections. (v) The welfare-relevant primitives
$(t,\alpha_1,\alpha_2)$ are identified from exogenous variation in opposite-side participation. Without such variation,
the reflection problem leaves $\alpha$ unidentified.
\end{abstract}

\section*{1. Model}
Sides $a\in\{1,2\}$ each have a unit mass of users uniform on $[0,1]$. Platform $1$ is at $0$ and platform $2$ at $1$. A side-$a$ user
at $x$ on platform $k$ gets $v+\alpha_aN_k^{-a}(\phi)-p_k^a-t\,d(x,k)$, where
$N_k^{-a}(\phi)=n_k^{-a}+\phi\,n_{-k}^{-a}$ is the accessible opposite side. Compatibility costs $\kappa(\phi)$, which is
increasing. With platform owners counted at unit weight, total welfare at allocation $(n^1,n^2)$, where $n^a$ is side $a$'s mass on
platform 1 (sorted by location), is
\[
 W(n;\phi)=2v+\sum_a\alpha_a\Bigl[n^an^{-a}+(1-n^a)(1-n^{-a})+\phi\bigl(n^a(1-n^{-a})+(1-n^a)n^{-a}\bigr)\Bigr]
 -\tfrac t2\sum_a\bigl[(n^a)^2+(1-n^a)^2\bigr]-\kappa(\phi).
\]

\begin{lemma}\label{lem:fb}
$\sup_nW(n;\phi)\le W^{\rm FB}:=2v+\alpha_1+\alpha_2-t/2-\kappa(\phi)$, with equality iff $\phi=1$ and $n^1=n^2=1/2$.
\end{lemma}
\begin{proof}
The bracket is at most $1$, with equality iff $\phi=1$ or the allocation is tipped. The transport term is minimised at
$n^a=1/2$, with value $t/4$ per side. At $\phi<1$ the two maxima are attained at different allocations.
\end{proof}

\section*{2. Equilibria and the threshold}
\begin{proposition}\label{prop:eq}
At $\phi=1$, each user's network benefit is $\alpha_a$ whatever platform it joins. At symmetric fees, the unique
participation equilibrium is $n^a=1/2$, so tipping cannot occur. At $\phi=0$ with symmetric fees, the split $n^a=1/2$ is
an equilibrium. It is the unique one if $t>\sqrt{\alpha_1\alpha_2}$ (more precisely, if the best-response map is a
contraction, for which $\max(\alpha_1,\alpha_2)<t$ suffices), and tipped allocations are equilibria if
$\alpha_a\ge t$ for both sides.
\end{proposition}
\begin{proof}
At $\phi=1$, the indifferent user on side $a$ solves $-tx=-t(1-x)$. At $\phi=0$, with symmetric fees, the indifferent
location is $x^a=\tfrac12+\tfrac{\alpha_a}{2t}(2n^{-a}-1)$. This map is a contraction when $\alpha_a<t$. If $\alpha_a\ge t$,
then $n^{-a}=1$ gives $x^a\ge1$, so tipping is self-confirming.
\end{proof}
\begin{theorem}[Selection-robust threshold]\label{thm:thr}
Let $\mathcal E(0)$ be the set of participation allocations sustainable in equilibrium at $\phi=0$, under any pricing
equilibrium selection. Then $W(e_1;1)\ge W(e_0;0)$ for all $e_0\in\mathcal E(0)$ iff
\[
 \kappa(1)-\kappa(0)\le\alpha_1+\alpha_2-\tfrac t2-\max_{e\in\mathcal E(0)}\bigl[W(e;0)+\kappa(0)-2v\bigr].
\]
For the split equilibrium the right side is $(\alpha_1+\alpha_2)/2$; for a tipped equilibrium it is $t/2$.
\end{theorem}
\begin{proof}
By Proposition~\ref{prop:eq}, $e_1$ is the split, and Lemma~\ref{lem:fb} gives $W(e_1;1)=W^{\rm FB}$. Compute
$W(\text{split};0)=2v+(\alpha_1+\alpha_2)/2-t/2-\kappa(0)$ and $W(\text{tipped};0)=2v+\alpha_1+\alpha_2-t-\kappa(0)$, then
compare.
\end{proof}
\begin{proposition}[Partial interoperability can hurt]\label{prop:partial}
Suppose $\alpha_a\ge t$ and the selection picks tipping at $\phi=0$. Let $\phi\in(0,1)$ be small enough that tipping is no
longer self-confirming, i.e.\ $(1-\phi)\alpha_a<t$, so that the split is selected. Then
$W(\text{split};\phi)-W(\text{tipped};0)=-(1-\phi)(\alpha_1+\alpha_2)/2+t/2-[\kappa(\phi)-\kappa(0)]$. This is negative
whenever $(1-\phi)(\alpha_1+\alpha_2)>t$.
\end{proposition}
\begin{proof}
At the split, the bracket equals $(1+\phi)/2$; substitute.
\end{proof}
So interoperability mandates should be either complete or accompanied by an argument that tipping was not efficient.

\section*{3. Identification}
\begin{proposition}\label{prop:id}
From observed shares and fees in a symmetric single-market equilibrium, $\alpha_a$ is not identified. The demand equation
$n^a=\tfrac12+\tfrac{\alpha_a(2n^{-a}-1)-(p_1^a-p_2^a)}{2t}$ involves the endogenous $n^{-a}$, and is observationally
equivalent to a model with $\alpha_a=0$ and correlated fee shocks (reflection). With an exogenous shifter $z$ of side
$-a$ participation, excluded from side-$a$ utility (for example a side-$(-a)$ subsidy), $(t,\alpha_a)$ are identified by IV:
$t$ from own-fee variation, and $\alpha_a/t$ from the $z$-induced variation in $n^{-a}$.
\end{proposition}
\begin{proof}
The structural equation is linear in $(p_1^a-p_2^a,\,2n^{-a}-1)$, with coefficients $(-1/2t,\ \alpha_a/2t)$. Two
instruments, the own-fee shifter and $z$, satisfy the rank condition. With only fees varying, $n^{-a}$ moves collinearly
with them in symmetric equilibria, and the rank condition fails.
\end{proof}

\section*{4. Scope}
Multi-homing, entry, quality investment and asymmetric platforms break the closed forms. Theorem~\ref{thm:thr}'s
structure persists, however: full interoperability removes the network-size component from users' platform choice and
therefore removes tipping-based multiplicity. Dynamic attraction basins and investment incentives (the statement's
remaining items) require a dynamic model, which we leave open \cite{rt,arm,weyl}.

\begin{thebibliography}{9}
\bibitem{rt} J.-C. Rochet, J. Tirole, Platform competition in two-sided markets, \emph{JEEA} 1 (2003), 990--1029.
\bibitem{arm} M. Armstrong, Competition in two-sided markets, \emph{RAND J. Econ.} 37 (2006), 668--691.
\bibitem{weyl} E. G. Weyl, A price theory of multi-sided platforms, \emph{AER} 100 (2010), 1642--1672.
\end{thebibliography}
\end{document}
